\chapter{Developer Documentation}
\label{ch:impl}

This chapter provides a comprehensive technical description of the air quality forecasting platform's implementation. It covers the system architecture, technology choices, data processing pipelines, model implementations, database design, security mechanisms, and the web application layer.

\section{Architecture Overview}
\label{sec:architecture}

\subsection{Layered Architecture}

The application follows a layered architecture pattern that separates concerns into distinct tiers, each responsible for a well-defined subset of functionality. This separation enables independent development, testing, and maintenance of each layer while maintaining clear interfaces between them.

The system is organized into four primary layers:

\begin{enumerate}
    \item \textbf{Presentation Layer} -- The Streamlit-based web interface responsible for rendering interactive visualizations, handling user input, and managing page navigation. This layer communicates exclusively with the business logic layer and never directly accesses the database or model artifacts.
    
    \item \textbf{Business Logic Layer} -- Contains the core application logic including model training orchestration, forecast generation, user management, and data validation. This layer acts as a mediator between the presentation and data layers.
    
    \item \textbf{Data Access Layer} -- Provides an abstraction over the SQLite database through dedicated repository classes. All SQL queries are encapsulated within this layer, exposing only high-level CRUD operations to the business logic layer.
    
    \item \textbf{Infrastructure Layer} -- Encompasses external dependencies including the file system (for model serialization), the SQLite database engine, and third-party APIs for data acquisition.
\end{enumerate}

\subsection{Component Interaction Diagram}

The following describes the component interaction topology of the system:

\begin{verbatim}
+----------------------------------------------------------+
|              front_end/ (Presentation Layer)               |
|  +------------+  +-------------+  +-------------------+  |
|  | Dashboard  |  | Predict     |  | Data Explorer     |  |
|  | History    |  | HelpDesk    |  | Settings          |  |
|  +------+-----+  +------+------+  +---------+---------+  |
+---------|--------------|-----------------------|----------+
          |              |                       |
+---------v--------------v-----------------------v----------+
|              back_end/ (Business Logic Layer)              |
|  +----------------+  +---------------+  +-------------+  |
|  | forecasting.py |  | data_proc.py  |  | helpdesk    |  |
|  | (LSTM/XGBoost) |  | (validation)  |  | (Gemini)    |  |
|  +-------+--------+  +-------+-------+  +------+------+  |
+---------|-----------------------|----------------|---------+
          |                       |                |
+---------v-----------------------v----------------v---------+
|              db/ (Data Access Layer)                        |
|  +-----------------------------------------------------+  |
|  | database.py: users, settings, history, chat, tokens |  |
|  +-----------------------------------------------------+  |
+-----------------------------------------------------------+
          |                       |
+---------v-----------------------v--------------------------+
|              Infrastructure                                |
|  +-------------+  +----------------+  +---------------+   |
|  | SQLite DB   |  | models/        |  | Gemini API    |   |
|  | (db/)       |  | (.keras,       |  | (external)    |   |
|  |             |  |  .joblib)      |  |               |   |
|  +-------------+  +----------------+  +---------------+   |
+-----------------------------------------------------------+
\end{verbatim}

\subsection{Data Flow}

The primary data flow through the system follows a pipeline pattern:

\begin{enumerate}
    \item \textbf{Data Ingestion}: Raw CSV or API data enters the system through the presentation layer's file upload mechanism or scheduled data fetching routines.
    
    \item \textbf{Validation and Transformation}: The business logic layer validates the schema, handles missing values, performs temporal indexing, and engineers features appropriate for each model type.
    
    \item \textbf{Model Training}: Processed data flows into the model training pipeline where it is split into training and testing sets, normalized as required, and used to fit the selected forecasting model.
    
    \item \textbf{Forecast Generation}: Trained models generate predictions that flow back through the business logic layer, are stored in the database, and are rendered as interactive visualizations in the presentation layer.
    
    \item \textbf{Persistence}: Model artifacts, training metadata, and forecast results are persisted to the file system and database respectively, enabling reproducibility and historical comparison.
\end{enumerate}

The system employs an event-driven interaction model within the Streamlit framework, where user actions (button clicks, parameter changes) trigger re-execution of the relevant computation pipeline. Session state is used to cache intermediate results and prevent unnecessary recomputation.

\subsection{Module Organization}

The project follows a layered architecture with clear separation between frontend, backend, database, and model concerns:

\begin{verbatim}
project_root/
--- front_end/              # Presentation layer (Streamlit)
-   --- Home.py             # Entry point
-   --- ui.py               # Shared layout, theme, session mgmt
-   --- charts.py           # Plotly configuration
-   --- pages/
-       --- 1_Data_Explorer.py
-       --- 2_Predict_future_quality.py
-       --- 3_Dashboard.py
-       --- 4_History.py
-       --- 5_Settings.py
-       --- 6_HelpDesk.py
-   --- static/theme.css    # Minimal CSS theme
--- back_end/               # Business logic layer
-   --- forecasting.py      # Model loading & 168h forecast
-   --- data_processing.py  # Data validation & preprocessing
-   --- page_dashboard.py   # Dashboard chart logic
-   --- page_predict.py     # Forecast figure builder
-   --- page_helpdesk.py    # Chat UI logic
-   --- helpdesk_gemini.py  # Gemini API client
--- db/                     # Data access layer
-   --- database.py         # SQLite operations (users, history,
-                           #   settings, chat, auth tokens)
--- models/                 # ML model layer
-   --- lstm_multivariate_model.keras
-   --- lstm_multivariate_artifacts.joblib
-   --- xgboost_multivariate_artifacts.joblib
-   --- lstm_forecast.py    # Standalone LSTM evaluation
-   --- xgboost_forecast.py # Standalone XGBoost evaluation
-   --- notebooks/          # Training scripts
--- data/                   # Datasets & preprocessing
-   --- london_2024.csv     # Primary dataset (hourly)
--- tests/                  # Pytest test suite
\end{verbatim}

\section{Technology Stack}
\label{sec:techstack}

\subsection{Python Ecosystem}

Python 3.10+ was selected as the primary development language for several compelling reasons:

\begin{itemize}
    \item \textbf{Scientific Computing Ecosystem}: Python offers the most mature and comprehensive ecosystem for data science and machine learning, with libraries such as NumPy, pandas, scikit-learn, TensorFlow, and XGBoost providing production-ready implementations of statistical and machine learning algorithms.
    
    \item \textbf{Rapid Prototyping}: Python's dynamic typing and expressive syntax enable rapid iteration during the research and development phase, which is critical for a thesis project with time constraints.
    
    \item \textbf{Community and Documentation}: The extensive community support and documentation available for Python's data science stack significantly reduces development friction.
    
    \item \textbf{Type Hinting Support}: Python 3.10+ provides robust type hinting capabilities that, combined with static analysis tools like mypy, offer many benefits of static typing without sacrificing development speed.
\end{itemize}

\subsection{Streamlit Framework Justification}

Streamlit was chosen over alternatives such as Flask+React, Dash, or Gradio for the following reasons:

\begin{table}[h]
\centering
\begin{tabular}{|l|c|c|c|c|}
\hline
\textbf{Criterion} & \textbf{Streamlit} & \textbf{Dash} & \textbf{Flask+React} & \textbf{Gradio} \\
\hline
Development Speed & Excellent & Good & Poor & Excellent \\
\hline
Interactive Widgets & Built-in & Built-in & Manual & Built-in \\
\hline
Data Visualization & Native & Native & Manual & Limited \\
\hline
Customizability & Moderate & High & Full & Low \\
\hline
Python-only Stack & Yes & Yes & No & Yes \\
\hline
Session Management & Built-in & Manual & Manual & Limited \\
\hline
Multi-page Support & Native & Manual & Manual & Limited \\
\hline
\end{tabular}
\caption{Web framework comparison for data science applications}
\label{tab:framework-comparison}
\end{table}

Streamlit's declarative programming model, where the script re-executes top-to-bottom on each interaction, aligns naturally with the data pipeline paradigm. Its built-in caching mechanism (\texttt{@st.cache\_data}, \texttt{@st.cache\_resource}) provides efficient memoization of expensive computations without requiring external caching infrastructure.

\subsection{Key Libraries}

\begin{table}[h]
\centering
\begin{tabular}{|l|l|p{7cm}|}
\hline
\textbf{Library} & \textbf{Version} & \textbf{Purpose} \\
\hline
pandas & $\geq$2.0 & DataFrame operations, temporal indexing, data manipulation \\
\hline
plotly & $\geq$5.15 & Interactive time series visualizations with zoom, pan, and hover \\
\hline
tensorflow & $\geq$2.13 & LSTM neural network implementation and training \\
\hline
xgboost & $\geq$1.7 & Gradient boosted tree ensemble model \\
\hline
statsmodels & $\geq$0.14 & ARIMA model fitting and statistical tests \\
\hline
scikit-learn & $\geq$1.3 & Preprocessing (MinMaxScaler), metrics, train/test splitting \\
\hline
bcrypt & $\geq$4.0 & Secure password hashing with adaptive cost factor \\
\hline
sqlite3 & stdlib & Embedded relational database (Python standard library) \\
\hline
streamlit & $\geq$1.28 & Web application framework \\
\hline
numpy & $\geq$1.24 & Numerical array operations \\
\hline
\end{tabular}
\caption{Primary library dependencies and their roles}
\label{tab:libraries}
\end{table}

The choice of Plotly over Matplotlib for visualization is motivated by Plotly's native interactivity (zooming into specific time ranges is essential for air quality analysis) and its seamless integration with Streamlit via \texttt{st.plotly\_chart()}.


\section{Data Processing Pipeline}
\label{sec:datapipeline}

The data processing pipeline transforms raw air quality measurements into model-ready feature matrices. This section details each stage of the pipeline, from initial loading through to the final train/test split.

\subsection{Data Loading and Schema Validation}

Raw data enters the system as CSV files containing timestamped air quality measurements. The loading procedure enforces a strict schema validation protocol:


The validation process performs the following checks:
\begin{itemize}
    \item Presence of all required columns
    \item Type coercion to expected data types
    \item Range validation (e.g., PM2.5 values must be non-negative)
    \item Duplicate timestamp detection and removal
    \item Chronological ordering verification
\end{itemize}

\subsection{Temporal Indexing}

After schema validation, the timestamp column is converted to a \texttt{DatetimeIndex} and set as the DataFrame index. This enables pandas' powerful time series functionality:


The frequency inference step is critical for ARIMA modeling, which requires regularly-spaced observations. When the inferred frequency is \texttt{None} (indicating irregular spacing), the data is resampled to hourly resolution using mean aggregation.

\subsection{Missing Value Handling}

Air quality sensor data frequently contains gaps due to sensor malfunctions, maintenance periods, or communication failures. The pipeline employs a multi-strategy approach to handle missing values:

\begin{enumerate}
    \item \textbf{Short gaps} ($\leq$ 3 consecutive missing values): Linear interpolation is applied, as short gaps in air quality data typically exhibit smooth transitions.
    
    \item \textbf{Medium gaps} (4--24 consecutive missing values): Forward fill followed by backward fill is used, preserving the last known state while avoiding extrapolation artifacts.
    
    \item \textbf{Long gaps} ($>$ 24 consecutive missing values): These segments are flagged and optionally excluded from training data, as imputation over such long periods would introduce significant bias.
\end{enumerate}


\subsection{Feature Engineering}

Each model type requires different feature representations. The pipeline generates model-specific feature sets:

\subsubsection{ARIMA Features}

ARIMA operates on univariate time series, so feature engineering is minimal:
\begin{itemize}
    \item Selection of the target pollutant column
    \item Stationarity testing via the Augmented Dickey-Fuller (ADF) test
    \item Differencing order determination based on ADF p-values
    \item Optional log transformation for variance stabilization
\end{itemize}

\subsubsection{LSTM Features}

The LSTM model requires normalized, windowed input sequences:


\subsubsection{XGBoost Features}

XGBoost requires tabular features with explicit lag and rolling statistics:


\subsection{Normalization}

Normalization strategies differ by model:

\begin{table}[h]
\centering
\begin{tabular}{|l|l|p{6cm}|}
\hline
\textbf{Model} & \textbf{Normalization} & \textbf{Rationale} \\
\hline
ARIMA & None & ARIMA operates on raw values; stationarity is achieved through differencing \\
\hline
LSTM & Min-Max [0,1] & Neural networks converge faster with bounded inputs; sigmoid/tanh activations expect normalized ranges \\
\hline
XGBoost & None & Tree-based models are invariant to monotonic transformations of features \\
\hline
\end{tabular}
\caption{Normalization strategy by model type}
\label{tab:normalization}
\end{table}

For the LSTM model, the \texttt{MinMaxScaler} is fitted exclusively on training data and then applied to both training and test sets. The scaler object is serialized alongside the model to enable proper inverse transformation during inference:

$$x_{\text{scaled}} = \frac{x - x_{\min}}{x_{\max} - x_{\min}}$$

\subsection{Train/Test Splitting Strategy}

Time series data requires temporal splitting rather than random splitting to prevent data leakage. The platform implements a chronological split:


The default split ratio is 80\% training / 20\% testing. This ratio was chosen to balance having sufficient training data for model convergence while retaining enough test data for meaningful evaluation. The assertion guard ensures no temporal leakage occurs between the sets.


\section{Model Implementations}
\label{sec:models}

This section provides detailed technical descriptions of the three forecasting models implemented in the platform: ARIMA, LSTM, and XGBoost. Each subsection covers the mathematical formulation, implementation details, and forecasting strategy.

\subsection{ARIMA -- AutoRegressive Integrated Moving Average}
\label{subsec:arima}

\subsubsection{Mathematical Formulation}

The ARIMA model combines three components: autoregression (AR), differencing (I), and moving average (MA). An ARIMA$(p, d, q)$ model is defined as follows.

Let $\{Y_t\}$ be the original time series. After $d$-th order differencing, we obtain the stationary series:

$$W_t = \nabla^d Y_t = (1 - B)^d Y_t$$

where $B$ is the backshift operator defined as $B^k Y_t = Y_{t-k}$, and $\nabla = 1 - B$ is the difference operator.

The ARIMA$(p, d, q)$ model for the differenced series is:

$$\phi(B) W_t = \theta(B) \varepsilon_t$$

where:
\begin{itemize}
    \item $\phi(B) = 1 - \phi_1 B - \phi_2 B^2 - \cdots - \phi_p B^p$ is the AR polynomial of order $p$
    \item $\theta(B) = 1 + \theta_1 B + \theta_2 B^2 + \cdots + \theta_q B^q$ is the MA polynomial of order $q$
    \item $\varepsilon_t \sim \text{WN}(0, \sigma^2)$ is white noise
\end{itemize}

Expanding the full model equation:

$$W_t = \phi_1 W_{t-1} + \phi_2 W_{t-2} + \cdots + \phi_p W_{t-p} + \varepsilon_t + \theta_1 \varepsilon_{t-1} + \theta_2 \varepsilon_{t-2} + \cdots + \theta_q \varepsilon_{t-q}$$

For the model to be valid, the AR polynomial must have all roots outside the unit circle (stationarity condition), and the MA polynomial must have all roots outside the unit circle (invertibility condition).

\subsubsection{Order Selection via Information Criteria}

The selection of optimal $(p, d, q)$ orders is performed using information criteria. The platform implements a grid search over candidate orders, evaluating both the Akaike Information Criterion (AIC) and the Bayesian Information Criterion (BIC):

$$\text{AIC} = -2 \ln(\hat{L}) + 2k$$

$$\text{BIC} = -2 \ln(\hat{L}) + k \ln(n)$$

where $\hat{L}$ is the maximized likelihood, $k$ is the number of estimated parameters, and $n$ is the number of observations.

The differencing order $d$ is determined prior to the grid search using the Augmented Dickey-Fuller (ADF) test. The null hypothesis $H_0: \text{unit root exists}$ is tested iteratively:


The grid search iterates over all combinations of $p \in \{0, 1, \ldots, 5\}$ and $q \in \{0, 1, \ldots, 5\}$, fitting each model and recording the information criterion value. Models that fail to converge are silently skipped. The default criterion is AIC, which tends to select slightly more complex models than BIC but provides better out-of-sample forecast accuracy for short-horizon predictions.

\subsubsection{Fitting Procedure}

Model fitting is performed via Maximum Likelihood Estimation (MLE). The log-likelihood for a Gaussian ARIMA model is:

$$\ln L(\phi, \theta, \sigma^2 | W) = -\frac{n}{2} \ln(2\pi) - \frac{n}{2} \ln(\sigma^2) - \frac{1}{2\sigma^2} \sum_{t=1}^{n} \varepsilon_t^2$$

The \texttt{statsmodels} implementation uses a state-space representation and the Kalman filter for efficient likelihood evaluation:


\subsubsection{Forecasting}

ARIMA forecasting produces point predictions along with confidence intervals. For an $h$-step ahead forecast:

$$\hat{Y}_{T+h|T} = E[Y_{T+h} | Y_1, Y_2, \ldots, Y_T]$$

The prediction intervals are computed as:

$$\hat{Y}_{T+h|T} \pm z_{\alpha/2} \cdot \sigma_h$$

where $\sigma_h^2 = \sigma^2 \sum_{j=0}^{h-1} \psi_j^2$ and $\psi_j$ are the coefficients of the infinite MA representation.


\subsubsection{Model Serialization}

ARIMA models are serialized using Python's \texttt{pickle} module, storing both the fitted model object and metadata:



\subsection{LSTM -- Long Short-Term Memory Network}
\label{subsec:lstm}

\subsubsection{Network Architecture}

Long Short-Term Memory (LSTM) networks are a specialized form of Recurrent Neural Networks (RNNs) designed to capture long-range temporal dependencies while mitigating the vanishing gradient problem. The LSTM cell maintains a cell state $C_t$ that acts as a conveyor belt for information flow, regulated by three gating mechanisms.

The LSTM cell at time step $t$ computes the following:

\textbf{Forget Gate:} Determines what information to discard from the cell state:
$$f_t = \sigma(W_f \cdot [h_{t-1}, x_t] + b_f)$$

\textbf{Input Gate:} Determines what new information to store:
$$i_t = \sigma(W_i \cdot [h_{t-1}, x_t] + b_i)$$

\textbf{Candidate Cell State:}
$$\tilde{C}_t = \tanh(W_C \cdot [h_{t-1}, x_t] + b_C)$$

\textbf{Cell State Update:}
$$C_t = f_t \odot C_{t-1} + i_t \odot \tilde{C}_t$$

\textbf{Output Gate:}
$$o_t = \sigma(W_o \cdot [h_{t-1}, x_t] + b_o)$$

\textbf{Hidden State:}
$$h_t = o_t \odot \tanh(C_t)$$

where $\sigma$ denotes the sigmoid activation function, $\odot$ represents element-wise multiplication, $W_*$ are weight matrices, and $b_*$ are bias vectors.

The platform implements a stacked LSTM architecture with the following topology:


The architecture consists of:
\begin{itemize}
    \item \textbf{Input shape}: \texttt{(window\_size, n\_features)} -- a 3D tensor where each sample is a sequence of \texttt{window\_size} time steps, each with \texttt{n\_features} measurements.
    \item \textbf{First LSTM layer}: 64 units with \texttt{return\_sequences=True} to pass full sequence output to the next LSTM layer.
    \item \textbf{Dropout}: 20\% dropout for regularization between LSTM layers.
    \item \textbf{Second LSTM layer}: 32 units with \texttt{return\_sequences=False}, outputting only the final hidden state.
    \item \textbf{Dense output layer}: Produces the forecast values for the next time step.
\end{itemize}
